\section{Entropy comparisons for every change of measure}
\label{l:tilts-section}

The domain transfers use marginal information about one fixed law.
Here the hypothesis instead compares entropy under every change of
measure supported on one common event. We condition on a large
representation coefficient and show that the resulting law puts mass
on complementary-restriction events that are rare under uniform
measure. The entropy comparison then bounds the probability of that
coefficient event. The common event must be chosen before the
representation and its two test vectors; an estimate only for the
original law would not suffice after this conditioning.

Let $(\Omega,\mathcal F,\mathbb P)$ be a probability space, and let
$g:\Omega\to S_n$ be a measurable endpoint map. For a law $\sigma$
on this space and $X\subseteq[n]$, denote the law of $g|_X$ by
$\sigma_X$ and its uniform injection reference by $U_X$.
Relative entropy uses natural logarithms:
\[
 D(\sigma\|\mathbb P)=
 \int\log\!\left(\frac{d\sigma}{d\mathbb P}\right)d\sigma
\]
when $\sigma\ll\mathbb P$, and is $+\infty$ otherwise.

\begin{theorem}[Arbitrary-tilt entropy transfer]
\label{l:tilted-entropy}
Let $128$ divide $n$, let $\mathcal G\in\mathcal F$, and suppose
$a=\mathbb P(\mathcal G)>0$. Assume that \emph{every} probability
law $\sigma$ on $(\Omega,\mathcal F)$ supported on $\mathcal G$
($\sigma(\mathcal G)=1$)
satisfies
\begin{equation}
 D(\sigma\|\mathbb P)\ge
 \E_{X:\,|X|=n-n/128}D(\sigma_X\|U_X)-n^{-2}.
 \label{l:all-tilts-entropy-input}
\end{equation}
There is an absolute $C$ such that, for every irreducible unitary
representation $\rho$ of $S_n$ of degree $D$ and all unit vectors $u,z$,
\begin{equation}
 \mathbb P\{\mathcal G,
       |\langle z,\rho(g)u\rangle|\ge D^{-1/8}\}
 \le CD^{-1/256}.
 \label{l:tilted-coefficient-tail}
\end{equation}
If $\eta$ is the endpoint law conditioned on $\mathcal G$, then
\begin{equation}
 \|\widehat\eta(\rho)\|_{\op}
 \le D^{-1/8}+\frac C aD^{-1/256}.
 \label{l:tilted-operator-bound}
\end{equation}
In particular, when $a\ge1/2$, this is at most $C'D^{-1/256}$
for an absolute $C'$. The statement therefore applies uniformly to a
sequence with $\mathbb P(\mathcal G)=1-o(1)$.
\end{theorem}

\begin{proof}
Put $b=128$. The claim for $D=1$ follows by increasing $C$, so assume
$D>1$. Fix $u,z$ and let $A$ be the event on the left side of
\eqref{l:tilted-coefficient-tail}. There is nothing to prove if
$\mathbb P(A)=0$. Otherwise set $\sigma=\mathbb P(\,\cdot\mid A)$.
Then $\sigma$ is supported on the same $\mathcal G$ as in the hypothesis,
and
\begin{equation}
 D(\sigma\|\mathbb P)=\log\frac1{\mathbb P(A)}.
 \label{l:tilted-conditioning-entropy}
\end{equation}
The reference law here remains $\mathbb P$, not
$\mathbb P(\,\cdot\mid\mathcal G)$. Thus this identity controls
the unconditioned probability of $A$; the factor $a^{-1}$ enters
only when we pass to the endpoint law $\eta$ at the end.
Our goal is to lower-bound this entropy by a positive multiple of
$\log D$. We do so using tests visible from complementary restrictions.

Fix an ordered partition $(C_1,\ldots,C_b)$ of $[n]$ into equal blocks.
Let $X_i=[n]\setminus C_i$ and $K_i=\Sym(C_i)$. On $V_\rho$, let
$S_i$ project onto the sum of the $K_i$-isotypic spaces with type degree
at most $D^{1/b}$. The subgroups $K_i$ commute, as do their isotypic
projections. On each tensor type in the product-group restriction, the
product of its factor degrees is at most $D$. At least one factor is
therefore at most $D^{1/b}$, giving
\[
 \prod_{i=1}^b(I-S_i)=0.
\]
Consequently
\begin{equation}
 u=\sum_{i=1}^b u_i,\qquad
 u_i=S_i\prod_{j<i}(I-S_j)u,\qquad \|u_i\|\le1.
 \label{l:tilted-vector-split}
\end{equation}

Let $W_i=\operatorname{span}\{\rho(h)u_i:h\in K_i\}$.
Lemma~\ref{l:single-orbit} bounds its dimension by the sum of the
squared type degrees, independently of their multiplicities.
Lemma~\ref{l:degree-reciprocal-twenty} supplies the finite constant
\[
 C_{20}=\sup_{m\ge1}
        \sum_{\tau\in\widehat{S_m}}(\dim\tau)^{-20}.
\]
It follows that
\begin{equation}
 \dim W_i\le\sum_{f\le D^{1/b}}f^2
       \le C_{20}D^{22/b},
 \label{l:tilted-orbit-rank}
\end{equation}
where the sum counts irreducible types once, not their multiplicities.

Because $W_i$ is $K_i$-invariant, $\rho(g)W_i$ depends only on
$g|_{X_i}$. Let $E_i$ be the subset of injections on $X_i$ for which
\[
 \|\operatorname{Proj}_{\rho(g)W_i}z\|
       \ge D^{-1/8}/b.
\]
Schur averaging gives uniform mean square projection length
$\dim W_i/D$. Markov's inequality and \eqref{l:tilted-orbit-rank} imply
\begin{equation}
 U_{X_i}(E_i)
 \le C_{20}b^2D^{-1+22/b+1/4}
 =C_{20}b^2D^{-37/64}
 \le C_{20}b^2D^{-1/2}.
 \label{l:tilted-uniform-test}
\end{equation}
For each endpoint on $A$, the coefficient threshold and
\eqref{l:tilted-vector-split} force at least one of these tests:
\[
 D^{-1/8}\le|\langle z,\rho(g)u\rangle|
 \le\sum_i\|\operatorname{Proj}_{\rho(g)W_i}z\|.
\]
Thus, for this partition,
\begin{equation}
 \sum_{i=1}^b\sigma_{X_i}(E_i)\ge1.
 \label{l:tilted-test-mass}
\end{equation}

We have found complementary restrictions on which $\sigma$ assigns
substantial mass to events that are rare under the uniform law. To
convert that fact to entropy, put $p_i=\sigma_{X_i}(E_i)$ and
$q_i=U_{X_i}(E_i)$. Data processing for the binary test gives
\[
 D(\sigma_{X_i}\|U_{X_i})
 \ge p_i\log\frac1{q_i}-\log2
 \ge\tfrac12p_i\log D-\log(2C_{20}b^2).
\]
The first inequality follows by expanding binary relative entropy and
using that binary Shannon entropy is at most $\log2$; its term
$(1-p_i)\log(1/(1-q_i))$ is nonnegative. Empty tests have $p_i=q_i=0$
and satisfy the displayed bound as well. The second inequality uses
\eqref{l:tilted-uniform-test} and $0\le p_i\le1$.

Average over $i$, using \eqref{l:tilted-test-mass}, and then over uniform
ordered equal-block partitions. Each $X_i$ is uniform of size $n-n/b$,
so
\[
 \E_{|X|=n-n/b}D(\sigma_X\|U_X)
 \ge\frac{\log D}{2b}-\log(2C_{20}b^2).
\]
Now apply \eqref{l:all-tilts-entropy-input} to this particular $\sigma$.
Together with \eqref{l:tilted-conditioning-entropy}, it gives
\[
 \log\frac1{\mathbb P(A)}
 \ge\frac{\log D}{256}-\log(2C_{20}b^2)-n^{-2}.
\]
This proves \eqref{l:tilted-coefficient-tail}, for example with
$C=2eC_{20}b^2$.

Finally the absolute coefficient is at most one. Split its expectation
under $\mathbb P(\,\cdot\mid\mathcal G)$ at $D^{-1/8}$, and use
\eqref{l:tilted-coefficient-tail} for the upper part. This proves
\eqref{l:tilted-operator-bound} after taking the supremum over the unit
vectors. All partitions and tests in the proof were chosen after $u,z$
were fixed; the hypothesis covering every law on the common event is
what permits the rare-event conditioning at the first step.
\end{proof}
